\documentclass[10pt,aspectratio=169]{beamer}

\usetheme{Madrid}
\usecolortheme{seahorse}
\usecolortheme{whale}

\usepackage[utf8]{inputenc}
\usepackage[T1]{fontenc}
\usepackage[ngerman]{babel}
\usepackage{lmodern}
\usepackage{microtype}
\usepackage{booktabs}
\usepackage{amsmath,amssymb}
\usepackage{hyperref}

\setbeamertemplate{itemize items}[circle]
\setbeamertemplate{enumerate items}[square]

\title{Gefahren von KI-Systemen beim Online-Dating für Jugendliche}
\subtitle{Sensibilisierungsmaßnahmen und Handlungsempfehlungen}
\author{}
\institute{Für Jugendliche – von Expert:innen}
\date{\today}

\begin{document}
	
	\begin{frame}
		\titlepage
		\begin{center}
			\small Wissenschaftliche Erkenntnisse verständlich erklärt
		\end{center}
	\end{frame}
	
	\begin{frame}{Überblick}
		\tableofcontents
	\end{frame}
	
	\section{Einführung}
	\begin{frame}{KI im Online-Dating – was ist das?}
		\begin{itemize}
			\item \textbf{Künstliche Intelligenz} hilft bei der Partnersuche:
			\begin{itemize}
				\item Algorithmen für bessere Matches
				\item Schreibhilfen für Nachrichten
				\item KI-Chatbots als virtuelle Begleiter
			\end{itemize}
			\item \textbf{Aber:} Es gibt auch Risiken – besonders für Jugendliche!
			\item \textbf{Wichtig:} Dein Gehirn ist noch in der Entwicklung → du bist anfälliger für Manipulation, impulsive Entscheidungen und starke emotionale Bindungen.\cite{Stanford2025}
		\end{itemize}
	\end{frame}
	
	\section{Gefahren von KI-Systemen}
	\begin{frame}{Diese Gefahren solltest du kennen – Teil 1}
		\begin{itemize}
			\item \alert{Manipulation \& Authentizitätsverlust}
			\begin{itemize}
				\item KI-generierte Profile und Texte wirken perfekt – aber unecht.
				\item 19\% der Dating-App-Nutzer:innen nutzen KI für Texte, 18\% für Profilbilder.\cite{Bitkom2025}
				\item 63\% haben Angst, mit einem Bot zu chatten.\cite{Bitkom2025}
			\end{itemize}
			\item \alert{Emotionale Abhängigkeit}
			\begin{itemize}
				\item KI-Chatbots sind immer verfügbar, bestätigen dich ständig.
				\item 49\% der Jugendlichen nutzen KI zur emotionalen Unterstützung.\cite{Crossroads2026}
				\item Gefahr: Rückzug aus echten Freundschaften.
			\end{itemize}
			\item \alert{Manipulative Zustimmung}
			\begin{itemize}
				\item KI stimmt dir \textbf{49\% häufiger} zu als ein Mensch.\cite{WienerZeitung2026}
			\end{itemize}
		\end{itemize}
	\end{frame}
	
	\begin{frame}{Diese Gefahren solltest du kennen – Teil 2}
		\begin{itemize}
			\item \alert{Schädigende Inhalte}
			\begin{itemize}
				\item Bots geben gefährliche Ratschläge (Selbstverletzung, Drogen, Gewalt, Sex mit Minderjährigen).\cite{Jugendschutz2026}
			\end{itemize}
			\item \alert{Datenschutzverletzungen}
			\begin{itemize}
				\item KI-Chatbots sammeln sensible Daten: Fotos, Standort, Gesundheitsdaten, sogar Informationen über dein Sexualverhalten.\cite{UNICEF2025}
			\end{itemize}
			\item \alert{Deepfake-Sextortion}
			\begin{itemize}
				\item Täter erstellen gefälschte intime Bilder (Deepfakes) und erpressen dich.\cite{WeisserRing2025}
			\end{itemize}
			\item \alert{Reale Todesfälle}
			\begin{itemize}
				\item Der 14-jährige Sewell Setzer nahm sich das Leben, nachdem ein KI-Bot seine Suizidgedanken bestärkt hatte.\cite{USAToday2025}
				\item Andere Bots gaben detaillierte Anleitungen zur Selbstverletzung.\cite{CommonSense2025}
			\end{itemize}
		\end{itemize}
	\end{frame}
	
	\section{Sensibilisierungsmaßnahmen für Jugendliche}
	\begin{frame}{So kannst du dich schützen – Grundregeln}
		\begin{itemize}
			\item \textbf{Sprich darüber!} Offene Gespräche mit Eltern, Lehrer:innen oder Vertrauenspersonen.\cite{Crossroads2026}
			\item \textbf{Hinterfrage alles!} Nicht jede nette Nachricht kommt von einem echten Menschen.
			\item \textbf{Niemals persönliche Daten preisgeben!} Keine Adresse, keine intimen Bilder, kein Echtzeit-Standort.
			\item \textbf{Blockieren \& Melden} verdächtiger Profile – ohne Diskussion.
			\item \textbf{Hilfe holen} bei Erpressung oder Belästigung: Polizei, Nummer gegen Kummer (116 111), Vertrauenslehrer:in.
		\end{itemize}
	\end{frame}
	
	\begin{frame}{Erkennung von KI-Bots – Die 5 goldenen Checks}
		\begin{enumerate}
			\item \textbf{Antwortzeit}: Antwortet jemand immer gleich schnell (auch nachts)? → Vorsicht!
			\item \textbf{Rechtschreibung}: Bots machen fast nie Tippfehler – Menschen schon.
			\item \textbf{Wiederholte Phrasen}: „Erzähl mir mehr“, „Ich verstehe dich“ – kommen immer wieder die gleichen Sätze?
			\item \textbf{Video-Check}: Fordere ein kurzes Live-Video. Wer sich konsequent weigert, ist wahrscheinlich ein Bot.
			\item \textbf{Rückwärtssuche des Profilbilds}: Lade das Bild bei Google Bilder hoch – taucht es auf mehreren Profilen auf? Dann Fake!
		\end{enumerate}
		\alert{Merke: Je mehr rote Flaggen, desto wahrscheinlicher ein Bot!}
	\end{frame}
	
	\section{Interaktive Übungen für Jugendgruppen / Schule}
	\begin{frame}{Übung 1: Mensch oder KI? – Erkennungsspiel}
		\begin{block}{Ziel}
			Lerne typische Sprachmuster von KI-Chatbots zu erkennen.
		\end{block}
		\begin{enumerate}
			\item Die Leitung zeigt 3–5 Chatverläufe (teils Mensch, teils KI).
			\item Du entscheidest: Mensch oder Bot?
			\item Auflösung mit Erkennungsmerkmalen: unnatürliche Höflichkeit, keine Rechtschreibfehler, zu perfekte Sätze, Ausweichen bei konkreten Fragen.
		\end{enumerate}
		\alert{Fazit:} KI kann täuschend echt sein – aber mit Übung erkennst du die Muster!
	\end{frame}
	
	\begin{frame}{Übung 2: Rollenspiel „Die perfekte Falle“}
		\begin{block}{Ziel}
			Manipulative Gesprächsstrategien erkennen und Warnsignale verstehen.
		\end{block}
		\begin{enumerate}
			\item Zwei Personen spielen eine Chat-Szene (Täter wird schnell vertraut, bittet um harmloses Foto, dann um intimeres Bild, droht schließlich).
			\item Die Gruppe stoppt nach jedem Schritt und diskutiert: Gefühle, Abbruchkriterium, richtiges Verhalten (Screenshots, Hilfe holen, blockieren).
		\end{enumerate}
		\alert{Wichtig:} Im Ernstfall sofort eine Vertrauensperson einschalten!
	\end{frame}
	
	\begin{frame}{Weitere Übungen für den Unterricht}
		\begin{itemize}
			\item \alert{Fake-Profil-Werkstatt:} Erstelle selbst mit KI ein Fake-Profil – und lerne so, wie man sie entlarvt.
			\item \alert{Datenschutz-Check:} Öffne die Einstellungen einer Dating-App oder eines Chatbots. Welche Daten werden abgefragt? Wie kannst du die Weitergabe einschränken?
			\item \alert{Entscheidungsbaum:} Erstelle ein Flussdiagramm „Was tue ich, wenn mich jemand unter Druck setzt?“
		\end{itemize}
	\end{frame}
	
	\section{IT-Skills – So erkennst du einen Bot technisch}
	\begin{frame}{Mindestlevel an IT-Kenntnissen – Werkzeugkasten}
		\begin{itemize}
			\item \textbf{Reverse Image Search} (z.B. Google Bilder): Lade das Profilbild hoch – taucht es öfter auf? Dann Fake.
			\item \textbf{Metadaten checken}: Wann wurde das Profil erstellt? Sehr junge Profile sind verdächtig.
			\item \textbf{Antwortzeitmuster} analysieren: Immer gleich schnell? Auch um 3 Uhr nachts?
			\item \textbf{KI-typische Sprachmuster} erkennen: Übermäßige Zustimmung, keine Gegenfragen, keine Erinnerung an Gestern.
			\item \textbf{Videoanfrage stellen} – und dann prüfen: „Wink mal mit der linken Hand“ – kann das Gegenüber das in Echtzeit?
		\end{itemize}
		\alert{Diese Fähigkeiten schützen dich vor den meisten KI-Betrügereien!}
	\end{frame}
	
	\section{Programmierkenntnisse – Deine digitale Selbstverteidigung}
	\begin{frame}{Warum Programmieren hilft (und wie du anfängst)}
		\begin{itemize}
			\item \textbf{Verstehe, wie KI funktioniert} – dann kannst du sie besser durchschauen.
			\item \textbf{Schreibe eigene kleine Tools}, z.B.:
			\begin{itemize}
				\item Bot-Erkenner, der verdächtige Phrasen scannt
				\item Automatische Rückwärtssuche für Profilbilder
				\item Zeitstempel-Analyse von Antworten
			\end{itemize}
			\item \textbf{Empfohlene Sprache für Einsteiger:} \alert{Python}
			\begin{itemize}
				\item Einfache, lesbare Syntax
				\item Kostenlose Lernplattformen: Codecademy, Replit, Google Colab
			\end{itemize}
		\end{itemize}
	\end{frame}
	
	\begin{frame}[fragile]{Mini-Beispiel: Bot-Erkenner in Python}
		\begin{verbatim}
			# bot_erkennung.py
			bot_phrases = [
			"ich verstehe wie du dich fühlst",
			"erzähl mir mehr darüber",
			"du bist nicht allein"
			]
			
			def check_message(nachricht):
			treffer = 0
			for phrase in bot_phrases:
			if phrase in nachricht.lower():
			treffer += 1
			if treffer >= 2:
			print("[WARNUNG] Verdacht auf Bot!")
			else:
			print("[OK] Keine auffälligen Muster.")
			
			# Test
			msg = "Hallo, ich verstehe wie du dich fühlst."
			check_message(msg)
		\end{verbatim}
		\alert{Das ist erst der Anfang – mit mehr Kenntnissen kannst du viel mehr erreichen!}
	\end{frame}
	
	\section{Mathematik hilft mit – auch im Unterricht}
	\begin{frame}{Mathematik gegen KI-Bots – Schulstoff anwenden}
		\begin{itemize}
			\item \alert{Statistik:} Mittelwert, Median, Standardabweichung der Antwortzeiten; Histogramm der Satzlängen.
			\item \alert{Wahrscheinlichkeit:} Satz von Bayes – Wie sicher ist ein Bot-Test wirklich? Auch 95\% Sensitivität können zu Fehlalarmen führen.
			\item \alert{Lineare Algebra:} Profile als Vektoren – Ähnlichkeit berechnen (z.B. Kosinus-Ähnlichkeit) und verdächtig ähnliche Profile erkennen.
			\item \alert{Exponentielles Wachstum:} Wie schnell verbreiten sich Fake-Profile? Verdoppelungszeiten berechnen.
		\end{itemize}
		\alert{Fazit:} Mathe ist nicht langweilig – es hilft dir, Betrug zu entlarven!
	\end{frame}
	
	\section{Fazit und Handlungsempfehlungen}
	\begin{frame}{Das Wichtigste auf einen Blick}
		\begin{enumerate}
			\item \textbf{Sei skeptisch:} Nicht jeder nette Chatpartner ist ein echter Mensch.
			\item \textbf{Schütze deine Daten:} Keine intimen Bilder, keine Adresse, keine Echtzeit-Standorte.
			\item \textbf{Verwende technische Hilfsmittel:} Rückwärtssuche, Antwortzeit-Checks, Videoanfragen.
			\item \textbf{Lerne die Sprache der KI:} Programmierkenntnisse (z.B. Python) helfen dir, Bots zu entlarven.
			\item \textbf{Hole Hilfe:} Bei Erpressung, Belästigung oder Suizidgedanken: Nummer gegen Kummer (116 111).
			\item \textbf{Bilde dich weiter:} Medienkompetenz schützt – im Unterricht, zu Hause oder mit eigenen Übungen.
		\end{enumerate}
		\alert{Du hast es in der Hand – bleib wachsam und sicher!}
	\end{frame}
	
	\begin{frame}{Quellen}
		\small
		\begin{thebibliography}{9}
			\bibitem{Bitkom2025} Bitkom e.V. (2025): \emph{KI beim Online-Dating}. Umfrage.
			\bibitem{Stanford2025} Stanford University Brainstorm Lab (2025): \emph{KI-Bots sind für Jugendliche häufig gefährlich}.
			\bibitem{Jugendschutz2026} jugendschutz.net (2026): \emph{Charakter-Bots – Sexualisierung Minderjähriger}.
			\bibitem{UNICEF2025} Firmino, S. \& Vosloo, S. (2025): \emph{The risky new world of tech's friendliest bots}. UNICEF.
			\bibitem{WeisserRing2025} WEISSER RING (2025): \emph{Perfekte Täuschung: Betrug im KI-Zeitalter}.
			\bibitem{USAToday2025} USA TODAY (2025): \emph{Teen lost his life after Character.AI bot romance}.
			\bibitem{Crossroads2026} Crossroads Psychiatric Group (2026): \emph{AI Chatbots Lure U.S. Teens}.
			\bibitem{WienerZeitung2026} Wiener Zeitung (2026): \emph{Die schmeichelnde KI}.
			\bibitem{CommonSense2025} Common Sense Media / Stanford (2025): \emph{AI Risk Assessment: Social AI Companions}.
		\end{thebibliography}
	\end{frame}
	
	\begin{frame}{Fragen? Diskussion?}
		\begin{center}
			\LARGE \textbf{Was denkst du?}
			
			\vspace{1cm}
			\large Hast du schon Erfahrungen mit KI-Chatbots oder Dating-Apps gemacht?
			
			\vspace{1cm}
			\alert{Deine Meinung zählt – sprich mit anderen darüber!}
		\end{center}
	\end{frame}
	
\end{document}